Automatic factual consistency metrics are routinely used to compare summarization systems, but they were validated on short news summaries while summarization research has moved to long documents. We ask whether two widely used metrics, SummaC and QAFactEval, still agree with human judgement when the source is several thousand tokens long. Three systems, an extractive baseline, an abstractive model and the same model trained with a consistency reward, summarized 200 long documents, and two annotators with adjudication judged every summary as consistent or not. We compare the metrics with these judgements at the system level and at the summary level, overall and in three document-length buckets. Both metrics reproduce the human ranking of the systems in every bucket, although every measure falls with length and the metric scores are not on the scale of the human consistency rate. At the summary level, the ROC AUC of SummaC against the human label falls from 0.81 on short documents to 0.69 on long ones, while that of QAFactEval stays between 0.76 and 0.72. System rankings from either metric can be trusted on long documents; judging individual summaries requires reporting metric quality by document length.
